\documentclass[11pt,a4paper]{article}
\usepackage[margin=1.9cm]{geometry}
\usepackage{fontspec}
\usepackage{polyglossia}
\setdefaultlanguage{greek}
\setotherlanguage{english}
\setmainfont{DejaVu Serif}[ItalicFont={DejaVu Serif},ItalicFeatures={FakeSlant=0.2},BoldItalicFont={DejaVu Serif Bold},BoldItalicFeatures={FakeSlant=0.2}]
\newfontfamily\greekfont{DejaVu Serif}[ItalicFont={DejaVu Serif},ItalicFeatures={FakeSlant=0.2},BoldItalicFont={DejaVu Serif Bold},BoldItalicFeatures={FakeSlant=0.2}]
\usepackage{amsmath,amssymb,amsthm}
\usepackage{enumitem}
\usepackage{fancyhdr}
\newcommand{\horrule}[1]{\rule{\linewidth}{#1}}
\newcommand{\eps}{\varepsilon}
\theoremstyle{definition}
\newtheorem{question}{Άσκηση}
\setlist[enumerate]{label=(\alph*),leftmargin=*,itemsep=3pt,topsep=5pt}
\setlength{\parindent}{0pt}
\setlength{\parskip}{5pt}
\setlength{\emergencystretch}{3em}
\pagestyle{fancy}
\fancyhf{}
\fancyfoot[L]{\small ΗΥ280 (2026) - Τσουρακάκης}
\fancyfoot[R]{\small\thepage}
\renewcommand{\headrulewidth}{0pt}

\begin{document}

\begin{center}
\horrule{0.4pt}\\[0.15cm]
{\large\textbf{ΗΥ 280 \textemdash{} Θεωρία Υπολογισμού}}\\[0.1cm]
{Ασκήσεις 7,8 \textemdash{} Κανονικές Γλώσσες, Κανονικές Εκφράσεις, ε-NFAs, DFAs}\\[0.1cm]
{Χαράλαμπος Ε. Τσουρακάκης}\\
\horrule{0.4pt}
\end{center}

\begin{center}\fbox{\begin{minipage}{0.93\linewidth}
\textbf{Οδηγίες για τις ασκήσεις}
\begin{itemize}[leftmargin=*,itemsep=6pt]
\item \textbf{Ασκήσεις 1-6, 7a, 8:} Αποτελούν μέρος του παραδοτέου. Μόνο το (β) ερώτημα της άσκησης 7 είναι προαιρετικό. 
\item Κάποιες από τις ασκήσεις θα συζητηθούν εκτενώς στο φροντιστήριο της Παρασκευής, μην το χάσετε! 
\end{itemize}
\end{minipage}}\end{center}\medskip

\begin{question}[Βασικές ιδιότητες πράξεων γλωσσών]

Έστω $\Sigma$ ένα πεπερασμένο αλφάβητο
και $A,B,L\subseteq\Sigma^*$ αυθαίρετες γλώσσες.

Να αποδείξετε τις ακόλουθες ισότητες:

\begin{enumerate}

\item Η συνένωση με την κενή γλώσσα δίνει
πάντοτε την κενή γλώσσα:
\[
    \boxed{
    \varnothing\cdot L
    =L\cdot\varnothing
    =\varnothing.
    }
\]

\item Η συνένωση με τη γλώσσα
που περιέχει μόνο την κενή λέξη
δεν αλλάζει τη γλώσσα:
\[
    \boxed{
    \{\eps\}\cdot L
    =L\cdot\{\eps\}
    =L.
    }
\]

\item Για οποιεσδήποτε γλώσσες $A,B$,
ισχύει:
\[
    \boxed{
    (A\cup B)^*=(A^*\cdot B^*)^*.
    }
\]

\textit{Υπόδειξη:}
Αποδείξτε τις δύο εγκλείσεις.
Χρησιμοποιήστε το γεγονός ότι
$\eps\in A^*, B^*$.

\end{enumerate}

\end{question}

\begin{question}[Περιγραφή γλωσσών]
Για κάθε μία από τις παρακάτω κανονικές εκφράσεις πάνω στο
$\Sigma=\{0,1\}$, περιγράψτε με λόγια τη γλώσσα που ορίζει.
\begin{enumerate}
\item $(0^*1^*)^*0$.
\item $(01^*)^*0$.
\item $(00+11+(01+10)(00+11)^*(01+10))^*$.
\end{enumerate}
\end{question}

\begin{question}[Απλοποίηση κανονικών εκφράσεων]
Απλοποιήστε τις ακόλουθες κανονικές εκφράσεις,
αιτιολογώντας τις απαντήσεις σας.
\begin{enumerate}
\item $(00)^*0+(00)^*$.
\item $(0+1)(\eps+00)^+ +(0+1)$.
\item $(0+\eps)0^*1$.
\end{enumerate}
\end{question}

\begin{question}[Δυαδικές αναπαραστάσεις δυνάμεων του 4]
Να βρείτε μια κανονική έκφραση για τη γλώσσα των
δυαδικών αναπαραστάσεων των αριθμών $4^n$, $n\ge 0$,
χωρίς αρχικά μηδενικά.
\end{question}

\begin{question}[Κατασκευή κανονικών εκφράσεων]
Για κάθε μία από τις παρακάτω γλώσσες πάνω στο
$\Sigma=\{0,1\}$, κατασκευάστε μία κανονική έκφραση.
\begin{enumerate}
\item Όλες οι συμβολοσειρές των οποίων το πέμπτο
σύμβολο από τα δεξιά (δηλαδή διαβάζοντας από τα δεξιά προς τα αριστερά) είναι $0$.
\item Όλες οι συμβολοσειρές που περιέχουν το $000$
ή το $111$ ως υποσυμβολοσειρά.
\item Όλες οι συμβολοσειρές που δεν περιέχουν
ούτε το $000$ ούτε το $111$ ως υποσυμβολοσειρά.
\end{enumerate}
\end{question}

\begin{question}[Ισοδυναμία κανονικών εκφράσεων]

Έστω το αλφάβητο $\Sigma=\{0,1\}$.
Γνωρίζουμε ότι η κανονική έκφραση
$(0+1)^*$ 
περιγράφει το σύνολο $\Sigma^*$ όλων
των δυαδικών λέξεων.
Να αποδείξετε ότι καθεμία από τις παρακάτω
κανονικές εκφράσεις είναι ισοδύναμη
με την $(0+1)^*$.

\begin{enumerate}

\item
$ 
    R_1=\eps+0(0+1)^*+1(1+0)^*.
$ 

\item
$R_2=0^*+0^*1(0+1)^*.$

\item
$ 
    R_3=((\eps+0)(\eps+1))^*.$

\item
$ R_4=0^*(10^*)^*.$

\end{enumerate}

\textit{Υπόδειξη:}
Για κάθε έκφραση $R_i$, αποδείξτε ότι $L(R_i)=\Sigma^*.$

Μπορείτε να χρησιμοποιήσετε τις δύο εγκλείσεις
\[
    L(R_i)\subseteq\Sigma^*,
    \qquad
    \Sigma^*\subseteq L(R_i).
\]

Η πρώτη έγκλειση είναι άμεση.
Για τη δεύτερη, θεωρήστε μία αυθαίρετη
δυαδική λέξη και βρείτε πώς μπορεί
να διασπαστεί σύμφωνα με τη δομή
της αντίστοιχης κανονικής έκφρασης.

\end{question}

\begin{question}[Αντιστροφή λέξεων και κλειστότητα κανονικών γλωσσών]

Έστω $\Sigma$ ένα πεπερασμένο αλφάβητο.

Υπενθυμίζουμε ότι η αντιστροφή $w^R$
μιας λέξης $w\in\Sigma^*$ ορίζεται
αναδρομικά ως εξής:
\[
w^R=
\begin{cases}
\eps,
    & \text{αν }w=\eps,\\[2mm]
x^R a,
    & \text{αν }w=ax,\quad
      a\in\Sigma,\ x\in\Sigma^*.
\end{cases}
\]

Για παράδειγμα,
\[
    (00101)^R=10100.
\]

Για κάθε γλώσσα $L\subseteq\Sigma^*$,
ορίζουμε την αντίστροφη γλώσσα
\[
    L^R=\{w^R\mid w\in L\}.
\]

\begin{enumerate}

\item Να αποδείξετε ότι, για κάθε γλώσσα
$L\subseteq\Sigma^*$,
\[
    \boxed{(L^*)^R=(L^R)^*.}
\]

\textit{Υπόδειξη:}
Χρησιμοποιήστε τον ορισμό
της επανάληψης Kleene και την ιδιότητα
\[
    (xy)^R=y^Rx^R.
\]

\item Να αποδείξετε ότι, αν η $L$
είναι κανονική γλώσσα,
τότε και η $L^R$ είναι κανονική.

\textit{Υπόδειξη:}
Ξεκινήστε από ένα DFA που αναγνωρίζει
τη $L$ και κατασκευάστε ένα
$\eps$-NFA για την $L^R$.

Αντιστρέψτε τη φορά όλων των
μεταβάσεων και προσέξτε τι πρέπει
να γίνει με τις αρχικές και
τις τελικές καταστάσεις.

Μπορείτε να χρησιμοποιήσετε
το αποτέλεσμα του (α).

\end{enumerate}

Μπορείτε να χρησιμοποιήσετε
χωρίς απόδειξη τις ακόλουθες ιδιότητες:
\[
\begin{aligned}
L^*L^*&=L^*,\\
(w^R)^R&=w,\\
(xy)^R&=y^Rx^R.
\end{aligned}
\]

\end{question}

\begin{question}[Κανονικές γλώσσες και διπλασιασμός συμβόλων]

Έστω $\Sigma$ ένα πεπερασμένο αλφάβητο. Ορίζουμε αναδρομικά τη συνάρτηση
\[
    \operatorname{stutter}:\Sigma^*\to\Sigma^*
\]
ως εξής:
\[
\operatorname{stutter}(w)=
\begin{cases}
\eps,
    & \text{αν } w=\eps,\\[2mm]
aa\cdot\operatorname{stutter}(x),
    & \text{αν } w=ax,
      \quad a\in\Sigma,\ x\in\Sigma^*.
\end{cases}
\]

Δηλαδή, η συνάρτηση $\operatorname{stutter}$
αντικαθιστά κάθε σύμβολο μιας λέξης
με δύο διαδοχικά αντίγραφά του.

Για παράδειγμα, αν $\Sigma=\{0,1\}$, τότε
\[
\begin{aligned}
\operatorname{stutter}(\eps)&=\eps,\\
\operatorname{stutter}(01)&=0011,\\
\operatorname{stutter}(101)&=110011.
\end{aligned}
\]

Για κάθε γλώσσα $L\subseteq\Sigma^*$,
ορίζουμε
\[
    \operatorname{stutter}(L)
    =
    \{\operatorname{stutter}(w)\mid w\in L\}.
\]

Έστω τώρα $L$ μία αυθαίρετη
κανονική γλώσσα. Να αποδείξετε ότι η γλώσσα
$    \operatorname{stutter}(L)$
είναι επίσης κανονική.

\textit{Υπόδειξη:}
Εφόσον η $L$ είναι κανονική,
υπάρχει DFA
\[
    M=(Q,\Sigma,\delta,q_0,F)
\]
που την αναγνωρίζει.

Κατασκευάστε ένα νέο αυτόματο $M'$ που,
για κάθε σύμβολο $a$,
προσομοιώνει μια μετάβαση
του $M$ διαβάζοντας το ζεύγος $aa$. Αποδείξτε με επαγωγή πως $L(M')=\operatorname{stutter}(L)$.
\end{question}

\end{document}
